\section*{Chapitre \ref{ch:dv:multi-asset-options} --- Options multi-actifs}

\begin{solution}{exo:dv:multi-asset-options:1}
$\sigma_B^2=0.25\times0.04+0.25\times0.09+2\times0.25\times0.5\times0.2\times0.3=0.0475$~: 21.8~\%, en dessous de la volatilité moyenne de
25~\%.
\end{solution}

\begin{solution}{exo:dv:multi-asset-options:2}
Un call best-of verse sur la plus forte performance, nettement au-dessus de la moyenne lorsque les actifs évoluent indépendamment et proche
de celle-ci lorsqu'ils évoluent ensemble~: la corrélation réduit son potentiel de hausse. Un call sur panier verse sur la moyenne, dont la volatilité croît avec
la corrélation~: moins de diversification, davantage de valeur d'option.
\end{solution}

\begin{solution}{exo:dv:multi-asset-options:3}
$\sum w_i^2\sigma_i^2=0.25(0.0625+0.1225)=0.04625$~; $(\sum w_i\sigma_i)^2=0.09$, donc le terme croisé vaut $0.04375$~;
$\bar\rho=(0.0576-0.04625)/0.04375=0.26$.
\end{solution}

\begin{solution}{exo:dv:multi-asset-options:4}
Soit $X$ le nombre d'unités de devise domestique par unité de devise étrangère. La valeur domestique de l'actif étranger, $SX$, est un actif négocié
en devise domestique, donc elle dérive de $r_d-q$ sous la mesure domestique. Le taux de change dérive de $r_d-r_f$. Selon Itô, la dérive de $S=(SX)/X$ est
$(r_d-q)-(r_d-r_f)+\sigma_X^2-\operatorname{cov}(\ln SX,\ln X)/dt$, et
$\operatorname{cov}(\ln SX,\ln X)/dt=\rho\sigma_S\sigma_X+\sigma_X^2$. La dérive vaut $r_f-q-\rho\sigma_S\sigma_X$, donc le forward quanto
est le forward étranger multiplié par $e^{-\rho\sigma_S\sigma_XT}$.
\end{solution}

\begin{solution}{exo:dv:multi-asset-options:5}
La note intègre un put worst-of vendu par le client (la jambe de capital à risque) et un coupon digital worst-of acheté par le client~;
le dealer est long du put et short du digital, et les deux perdent de la valeur pour lui à mesure que la corrélation monte. Des options sur indice, ou
la variance d'indice contre la variance des actions individuelles (une dispersion inversée), portent l'exposition inverse à la corrélation et compensent une partie
de cette exposition.
\end{solution}

\begin{solution}{exo:dv:multi-asset-options:6}
Environ $+28$~: les montants de couverture annulent l'exposition du premier ordre à la variance de chaque membre, et ce qui reste est un terme du second
ordre, faible comparé aux 270\,600 que produit un écart de corrélation de quinze points.
\end{solution}

\begin{solution}{exo:dv:multi-asset-options:7}
0.736, entre les prix à corrélation constante de 0.6 (0.707) et de 0.8 (0.748), et bien au-dessus du 0.673 obtenu à la
\omterm{def:dv:multi-asset-options:implied}{corrélation implicite} à la monnaie de l'indice, d'environ 0.4~: avec la \omterm{def:dv:multi-asset-options:local}{corrélation locale}, les actions évoluent ensemble exactement dans les scénarios de baisse
qui décident du digital.
\end{solution}

\begin{solution}{exo:dv:multi-asset-options:8}
C'est une prime, et non un carry~: le livre est short de corrélation et short de volatilité d'indice, et il perd lourdement dans un repli,
lorsque la corrélation tend vers un et que la volatilité de l'indice s'envole. Les travaux publiés montrent aussi que des frictions de négociation réalistes
rendent la prime difficile à capturer. L'écart moyen rémunère ce risque.
\end{solution}

\begin{solution}{pb:dv:multi-asset-options:1}
\textbf{1.} $\sum w_i^2\sigma_i^2=0.00430$ (43.0 points de variance)~; $C=0.0798$ (798).
\textbf{2.} 6.6~\% à corrélation nulle~; 29.0~\% à corrélation égale à un, soit la volatilité moyenne pondérée.
\textbf{3.} Diversification~: avec des corrélations inférieures à un, les mouvements des membres se compensent en partie.
\textbf{4.} 21.9~\%, 19.2~\% et 17.0~\%.
\textbf{5.} 19.3~\% dans les deux cas~: un smile d'indice plat.
\textbf{6.} 0.549, 0.408 et 0.308.
\textbf{7.} 0.559, à partir de la volatilité du strip de l'indice, de 22.1~\%.
\textbf{8.} Les options sur indice hors de la monnaie valorisent les membres qui baissent ensemble, ce qui correspond à une corrélation plus élevée.
\textbf{9.} $\rho(x)=0.384-3.62x+16.3x^2$, bornée à $[0,1]$.
\textbf{10.} 0.91 pour $x=-0.1$ et 0.19 pour $x=+0.1$.
\textbf{11.} $100\,000/(2\times22.12)=2\,261$ par point de variance.
\textbf{12.} $N\bigl(w_i^2(1-\bar\rho)+\bar\rho w_i\sum_jw_j\sigma_j/\sigma_i\bigr)$~: la sensibilité de la variance de l'indice à la variance de chaque membre
à la \omterm{def:dv:multi-asset-options:implied}{corrélation implicite}~; elles annulent l'exposition du premier ordre à la volatilité de chaque membre.
\textbf{13.} $2\,261\times0.15\times798=270\,600$.
\textbf{14.} 271\,000, 34\,000, de 215\,000 à 326\,000.
\textbf{15.} La corrélation moyenne réalisée, 0.409, moins 0.559~: $-0.15$ par unité de notionnel pour le côté long, $+0.15$ pour le
côté court.
\textbf{16.} Lorsque la corrélation se réalise au-dessus de l'implicite, typiquement lors d'un repli~; la perte vaut $N(\hat\rho-\bar\rho)C$ plus le
résidu lié aux variations de volatilité, et la corrélation peut aller jusqu'à un~: environ $2\,261\times0.441\times798\approx800\,000$ au titre de la seule corrélation
si les membres réalisent leurs volatilités implicites.
\textbf{17.} La volatilité implicite a été en moyenne supérieure à la volatilité réalisée~: 39.5~\% contre 32.5~\% pour le S\&P~500, 46.0~\% contre 35.5~\% pour le Dow Jones 30,
dans l'échantillon de Driessen, Maenhout et Vilkov.
\textbf{18.} Les investisseurs paient pour couvrir les pics de corrélation (ils sont acheteurs de puts sur indice), et le trade qui la vend exige de nombreuses
positions sur options et supporte un risque de krach~; les frictions limitent l'arbitrage.
\textbf{19.} \omterm{def:dv:multi-asset-options:implied}{Corrélation implicite} de 0.549 à 90, 0.408 à 100 et 0.308 à 110 (0.559 en termes de variance)~; le \omterm{def:dv:multi-asset-options:dispersion}{trade de dispersion} rapporte
270\,600 sur un \omterm{def:dv:variance-swaps-and-volatility-derivatives:veganot}{notionnel vega} de 100\,000 lorsque la corrélation réalisée est inférieure de quinze points à l'implicite (simulation~: 271\,000 $\pm$ 34\,000).
\textbf{20.} Le comportement conjoint de ses membres, surtout la manière dont leur corrélation évolue lorsque le marché baisse.
\end{solution}

\begin{solution}{iq:dv:multi-asset-options:1}
La valeur d'un call sur panier croît avec la corrélation (moins de diversification, plus de volatilité de la moyenne)~; un put worst-of et un call best-of
baissent avec elle (les extrêmes des performances se rapprochent de la moyenne à mesure que la corrélation augmente).
\iqlookfor{les trois signes et le mécanisme.}
\end{solution}

\begin{solution}{iq:dv:multi-asset-options:2}
La corrélation unique qui réconcilie la variance implicite de l'indice avec celles de ses membres~:
$(\sigma_I^2-\sum w_i^2\sigma_i^2)/(\text{terme croisé})$. Les skews d'indice sont plus marqués que les skews des actions individuelles, de sorte que la volatilité de l'indice aux strikes bas implique une corrélation plus élevée~: le
marché valorise des actions qui baissent ensemble.
\iqlookfor{la formule et le \omterm{def:dv:multi-asset-options:skew}{skew de corrélation}.}
\end{solution}

\begin{solution}{iq:dv:multi-asset-options:3}
Vente de variance d'indice (ou de straddles), achat de variance des membres en montants annulant l'exposition de chaque membre à sa volatilité, à la
\omterm{def:dv:multi-asset-options:implied}{corrélation implicite}. Le résidu est une position short de corrélation~: P\&L $\approx N(\bar\rho-\hat\rho)C$. Exposition aux pics de corrélation lors des replis,
à la volatilité des pondérations et des ratios de couverture, et aux coûts de transaction sur de nombreuses jambes.
\iqlookfor{les jambes, les pondérations et l'exposition résiduelle.}
\end{solution}

\begin{solution}{iq:dv:multi-asset-options:4}
Sous la mesure domestique, la dérive de l'actif étranger vaut $r_f-q-\rho\sigma_S\sigma_X$, avec $X$ exprimé en unités domestiques par unité étrangère (Girsanov
lors du changement de numéraire). Pour un indice libellé en yens et payé en dollars, $X$ vaut des dollars par yen~; si l'indice baisse quand le yen se renforce,
$\rho<0$, et le forward quanto se situe au-dessus du forward en yens.
\iqlookfor{la dérive et la convention de signe.}
\end{solution}

\begin{solution}{iq:dv:multi-asset-options:5}
La matrice n'est pas définie positive~: les estimations par paires issues de fenêtres ou de données différentes, ou les valeurs saisies à la main, sont incohérentes. Il faut la
projeter sur la matrice de corrélation la plus proche (en tronquant les valeurs propres négatives et en renormalisant la diagonale, ou par une méthode de projections alternées), ou
la construire à partir d'un modèle factoriel, et vérifier la variation des prix.
\iqlookfor{la définie-positivité et une réparation rigoureuse.}
\end{solution}

\begin{solution}{iq:dv:multi-asset-options:6}
Mesure~: la variation de valeur du livre pour un choc sur l'ensemble des corrélations deux à deux (et pour un choc de \omterm{def:dv:multi-asset-options:skew}{skew de corrélation}), par compartiment de maturité
et de la moneyness du plus mauvais performer. Couverture~: avec des options sur indice ou de la variance d'indice contre la variance des actions individuelles, qui portent
une exposition à la corrélation de signe opposé, en sachant que la couverture est imparfaite (la corrélation de l'indice n'est pas celle du panier). Réserve~: pour la
corrélation non couverte à un niveau de stress et pour le modèle (\omterm{def:dv:multi-asset-options:local}{corrélation locale} contre corrélation constante).
\iqlookfor{une sensibilité à la corrélation, une couverture par proxy et une réserve.}
\end{solution}
